\section{Introduction}
\label{sec:intro}

% ARC: exotics renaissance -> do multiquark states form often or rarely?
% -> Lebed's static estimate -> two gaps: dynamics, and rigor -> this paper.

The proliferation of experimentally established multiquark candidates ---
$\chic$~\cite{Choi:2003ue,LHCb:2020xds,CMS:2021znk}, $\Tcc$~\cite{LHCb:2021vvq,LHCb:2021auc},
the $P_c$ pentaquarks~\cite{LHCb:2015yax,LHCb:2019kea} --- has renewed the
question of how frequently quarks assemble into hadrons beyond $\qqbar$ and
$qqq$.  Lebed~\cite{Lebed:2016jvq} gave a strikingly simple estimate:
in an ideal gas of quarks and antiquarks, the nearest-neighbor distance
distribution~\cite{Chandrasekhar:1943ws}
\begin{equation}
w(r)\,dr = 4\pi r^2 n\, e^{-4\pi r^3 n/3}\, dr
\label{eq:nn}
\end{equation}
combined with single-gluon-exchange color factors implies that a quark is
attracted more strongly to its nearest fellow quark than to its nearest
antiquark tens of percent of the time.  Diquark formation, and by extension
multiquark-hadron formation, should be \emph{common}.

That estimate is static: quarks neither move, nor collide, nor annihilate,
nor bind.  Its counting is also deliberately schematic.  This paper closes
both gaps.

First (Secs.~\ref{sec:model}--\ref{sec:classical-results}), we promote
the model to a dynamical simulation: a molecular-dynamics evolution of
an expanding quark--antiquark plasma with discrete color charges and
Casimir-weighted screened interactions, in which binding is emergent
(the Hubble redshift supplies the dissipation), annihilation proceeds
through tunable-timescale photon and gluon channels, and hadron species
are identified as persistent color-singlet clusters.  Production
ensembles ($\sim1100$ events up to $N=10^4$ quarks) yield three
headline results: the species ladder is strongly non-geometric
(meson$\to$baryon costs $\sim65\times$ but baryon$\to$tetraquark only
$\sim5\times$); manifestly flavor-exotic tetraquarks form at
$0.3$/event under strong coupling while double-charm exotics are
doubly suppressed (seeded $cc$ diquarks melt above their
$\sim36\MeV$ binding); and the fraction of baryons assembled through
a persistent diquark intermediate is $0.91$--$0.95\pm0.02$ across the
entire coupling--screening plane --- dynamics \emph{amplifies} the
static estimate by more than a factor of three.

Second (Secs.~\ref{sec:quantum-setup}--\ref{sec:protocols}), we ask what it
would take to answer the same question \emph{rigorously}, from QCD itself.
Real-time hadronization is precisely the regime where Euclidean lattice
methods fail and quantum simulation is expected to
excel~\cite{Bauer:2022hpo,JordanLeePreskill:2012}.  We take an ideal
fault-tolerant quantum computer executing $3{+}1$D lattice QCD as given, and
construct the measurement layer: gauge-invariant,
renormalization-group-safe definitions of ``number of mesons, baryons,
tetraquarks, pentaquarks'' in the out-state of a fireball quench, and the
matrix elements that extract their ratios
(Sec.~\ref{sec:hierarchy}).  The central observation is that species
numbers built from (i) integer-quantized conserved charges on isolated
color-singlet clusters and (ii) spectral data of those clusters are
RG-invariant, while any definition that reaches inside a hadron is
irreducibly scheme-dependent and must be labeled as such
(Sec.~\ref{sec:tier4}).  A second observation defuses the finite-time
objection: the flagship exotics are narrow
($c\tau_{\chic}\!\approx\!165\fm$, $c\tau_{\Tcc}\!\approx\!4000\fm$), so
there is a parametric window $t_{\rm fo}\ll t_f \ll 1/\Gamma$ in which they
are countable clusters (Sec.~\ref{sec:resonances}).

Both halves speak to an existing phenomenological debate.  Statistical
hadronization \cite{Andronic:2017pug} populates species by mass and
quantum numbers at a universal chemical freeze-out temperature and is
agnostic about internal structure, while coalescence
\cite{Fries:2008hs} rewards compact wavefunctions and small constituent
counts; the ExHIC program \cite{ExHIC:2010gcb,ExHIC:2017smd} showed
that the \emph{ratio} of these two predictions for a given exotic is
itself a structure diagnostic (molecular states are coalescence-enhanced,
compact ones suppressed), a logic since applied to $\chic$ yields in
heavy-ion data \cite{Zhang:2020dwn,CMS:2021znk}.  Our classical model is
in effect a dynamical coalescence calculation with tunable formation and
survival timescales, and the quantum framework renders the same
observables field-theoretically well defined; both are confronted with
the phenomenological baselines in Sec.~\ref{sec:confrontation}.
